\documentclass{article}

\pagestyle{empty}

\begin{document}

Here is an example of some normal text. As you see, \LaTeX\ typesets text in fully justified mode by default. This works by adding spacing between words in order to make the text look nice. But sometimes, the system needs to hyphenate a word if it is too long to fit on a line and look good. As you can see, Donald Knuth put a lot of thought into the algorithms needed to make the text look nice! But you can either turn hyphenation off or seize control of the hyphenation of particular words if you are not satisfied with the results. \\

\begin{flushleft}
Here is an example of flush left text. As you can see, the text looks a bit different from normal text. Specifically, there is no more attempt to make the text fill the entire line by inserting extra spacing between words. Instead, the text is flush on the left margin only. The right margin is left ragged.\\

\end{flushleft}

\begin{flushright}
Here is an example of flush right text. As you can see, flush right behaves similarly to flush left. The only real difference is that the the text is flush with the right margin instead of the left margin.  Now, it is the left margin that is left ragged.\\

\end{flushright}

\begin{center}
And finally, this is an example of centered text. It's pretty obvious what this environment does without much in the way of explanation. \\

\end{center}

\end{document}




